% !TeX root = part8.tex
\documentclass[../final.tex]{subfiles}
\begin{document}
\ifSubfilesClassLoaded{\setcounter{chapter}{5}}{}

% Sources: Desktop/11.jpg, 22.jpg, 33.jpg.
\chapter{Фононы в твёрдом теле}
\rsubtitle{Модель Дебая, теплоёмкость и колебания атомов}

\section{Нормальные колебания решётки}
\subsection{Смещения атомов и гармоническое приближение}
Рассмотрим одноатомное твёрдое тело из $N$ атомов одинаковой массы $m$.
Положение $i$-го атома представим как
\begin{equation*}
 \mathbf r_i=\mathbf R_i+\mathbf u_i,
 \qquad\boldsymbol\chi=(u_{1x},u_{1y},u_{1z},\ldots,u_{Nz})^{\mathsf T},
\end{equation*}
где $\mathbf R_i$ --- положение равновесия, а $\mathbf u_i$ --- малое смещение.
У каждого атома три компоненты смещения, всего --- $3N$ координат $\chi_a$.
Разложим потенциальную энергию около равновесия:
\begin{align*}
 U(\boldsymbol\chi)&=U_{\mathrm{stat}}
 +\sum_a\left.\frac{\partial U}{\partial\chi_a}\right|_0\chi_a
 +\frac12\sum_{a,b}\left.\frac{\partial^2U}
 {\partial\chi_a\partial\chi_b}\right|_0\chi_a\chi_b+\cdots,\\
 \left.\frac{\partial U}{\partial\chi_a}\right|_0&=0,\qquad
 B_{ab}=\frac1m\left.\frac{\partial^2U}
 {\partial\chi_a\partial\chi_b}\right|_0.
\end{align*}
Линейный член исчезает, поскольку в равновесии силы равны нулю.
Гармоническое приближение сохраняет только квадратичный член:
\begin{equation*}
 U_2\equiv U-U_{\mathrm{stat}}
 =\frac m2\sum_{a,b}B_{ab}\chi_a\chi_b
 =\frac m2\boldsymbol\chi^{\mathsf T}B\boldsymbol\chi.
\end{equation*}
Матрица $B$ вещественна и симметрична: $B_{ab}=B_{ba}$.
Смешанные произведения $\chi_a\chi_b$ при $a\ne b$ описывают связь
между смещениями разных координат. Движения отдельных атомов поэтому
не являются независимыми осцилляторами.
С учётом кинетической энергии гамильтониан равен
\begin{equation*}
 H=U_{\mathrm{stat}}+\frac{\mathbf p^{\mathsf T}\mathbf p}{2m}
 +\frac m2\boldsymbol\chi^{\mathsf T}B\boldsymbol\chi.
\end{equation*}
Множитель $m/2$ можно включить в коэффициенты квадратичной формы.
Принятое здесь определение $B$ удобно тем, что её собственные значения
имеют размерность квадрата частоты.

\clearpage
\subsection{Почему внутренних колебаний $3N-6$}
Три перемещения центра масс и три поворота свободного конечного нелинейного тела
не меняют взаимных расстояний атомов. После их отделения остаётся
$M=3N-6$ внутренних мод; при $N\gg1$ имеем $M\simeq3N$.

\subsection{Матрица поворота и нормальные координаты}
Вещественная симметричная матрица $B$ имеет ортонормированный набор
собственных векторов $\mathbf e_\nu$:
\begin{equation*}
 B\mathbf e_\nu=\omega_\nu^2\mathbf e_\nu,\qquad
 \mathbf e_\nu^{\mathsf T}\mathbf e_\mu=\delta_{\nu\mu}.
\end{equation*}
Составим из этих векторов столбцы матрицы $Z$.
Это ортогональное преобразование базиса:
\begin{equation*}
 Z^{\mathsf T}Z=1,\qquad Z^{-1}=Z^{\mathsf T},\qquad
 Z^{\mathsf T}BZ=\Omega^2
 =\operatorname{diag}(\omega_1^2,\omega_2^2,\ldots).
\end{equation*}
Нормальные координаты --- проекции смещений на собственные векторы:
\begin{equation*}
 \mathbf Q=Z^{\mathsf T}\boldsymbol\chi,\qquad
 \boldsymbol\chi=Z\mathbf Q,\qquad
 Q_\nu=\sum_a Z_{a\nu}\chi_a.
\end{equation*}
Подстановка в потенциальную энергию даёт всю цепочку преобразований:
\begin{align*}
 U_2&=\frac m2\boldsymbol\chi^{\mathsf T}B\boldsymbol\chi
 =\frac m2(Z\mathbf Q)^{\mathsf T}B(Z\mathbf Q)\\
 &=\frac m2\mathbf Q^{\mathsf T}(Z^{\mathsf T}BZ)\mathbf Q
 =\frac m2\mathbf Q^{\mathsf T}\Omega^2\mathbf Q
 =\frac m2\sum_\nu\omega_\nu^2Q_\nu^2.
\end{align*}
Смешанные произведения нормальных координат исчезли.
Преобразуем тем же способом импульсы: $\boldsymbol\Pi=Z^{\mathsf T}\mathbf p$.
Ортогональность сохраняет кинетическую энергию,
$\mathbf p^{\mathsf T}\mathbf p=\boldsymbol\Pi^{\mathsf T}\boldsymbol\Pi$.
После отделения движения тела как целого
\begin{equation*}
 H_{\mathrm{int}}=U_{\mathrm{stat}}+\sum_{\nu=1}^{M}
 \left(\frac{\Pi_\nu^2}{2m}+\frac{m\omega_\nu^2Q_\nu^2}{2}\right),\qquad
 \ddot Q_\nu+\omega_\nu^2Q_\nu=0.
\end{equation*}
Каждая нормальная координата колеблется независимо с собственной частотой
$\omega_\nu$. Одна мода описывает согласованное движение многих атомов.

\subsection{Продольные и поперечные волны}
Нормальные колебания являются коллективными движениями атомов.
При продольной волне смещения направлены вдоль волнового вектора,
при поперечной --- перпендикулярно ему.

\begin{rbox}[unbreakable,left=14pt,right=14pt,top=8pt,bottom=6pt]{[]}
\centering
\captionsetup{font=footnotesize,hypcap=false,skip=5pt}
\begin{tikzpicture}[>=Stealth,font=\small,x=0.95cm,y=0.45cm]
 \draw[->] (0,0)--(11.6,0) node[right] {$\mathbf k$};
 \foreach \xx in {0.4,1.25,2.0,2.6,3.1,3.55,4.2,5.0,6.0,7.0,7.8,8.4,8.9,9.4,10.2} {
  \fill[p3] (\xx,0) circle (1.7pt);
 }
 \draw[<->,red3] (2.8,0.55)--(4.2,0.55);
 \node[above] at (7,0.55) {продольная волна};
 \draw[->] (0,-2.1)--(11.6,-2.1) node[right] {$\mathbf k$};
 \draw[p3,domain=0:11,samples=100] plot (\x,{-2.1+0.42*sin(65*\x)});
 \foreach \xx in {0.4,1.2,2,2.8,3.6,4.4,5.2,6,6.8,7.6,8.4,9.2,10,10.8} {
  \fill[p3] (\xx,{-2.1+0.42*sin(65*\xx)}) circle (1.7pt);
 }
 \draw[<->,red3] (1.3,-2.7)--(1.3,-1.5);
 \node[below] at (7,-2.7) {поперечная волна};
\end{tikzpicture}
\captionof{figure}{Продольная и поперечная акустические волны в решётке.}
\end{rbox}

\clearpage
\section{Квантовые осцилляторы и фононы}
\subsection{Статистическая сумма одной моды}
Температура в формулах измеряется в энергетических единицах:
$T=k_BT_{\mathrm K}$, $\beta=1/T$.
Энергетические уровни гармонического осциллятора равны
\begin{equation*}
 E_n=\hbar\omega\left(n+\frac12\right),\qquad n=0,1,2,\ldots.
\end{equation*}
Суммируя геометрическую прогрессию, получаем
\begin{align*}
 Z_\omega&=\sum_{n=0}^{\infty}e^{-\beta\hbar\omega(n+1/2)}\\
 &=e^{-\beta\hbar\omega/2}
 \left(1+e^{-\beta\hbar\omega}+e^{-2\beta\hbar\omega}+\cdots\right)
 =\frac{e^{-\beta\hbar\omega/2}}{1-e^{-\beta\hbar\omega}},\\
 \ln Z_\omega&=-\frac{\beta\hbar\omega}{2}
 -\ln(1-e^{-\beta\hbar\omega}).
\end{align*}
Средняя энергия моды равна
\begin{equation*}
 \overline E_\omega=-\frac{\partial\ln Z_\omega}{\partial\beta}
 =\frac{\hbar\omega}{2}+\frac{\hbar\omega}{e^{\hbar\omega/T}-1}.
\end{equation*}
Первое слагаемое --- энергия нулевых колебаний, второе --- тепловая энергия.
Квант возбуждения нормальной моды называется \rconcept{фононом}.

\subsection{Число фононов и сравнение с фотонами}
Число атомов $N$ фиксировано, а число фононов $N_{\mathrm{ph}}$ не сохраняется.
Для равновесных фононов $\mu_{\mathrm{ph}}=0$, поэтому
\begin{equation*}
 \overline n(\omega)=\frac1{e^{\hbar\omega/T}-1},\qquad
 N_{\mathrm{ph}}=\sum_\nu\overline n(\omega_\nu).
\end{equation*}
У фотонов $\omega=ck$ и две поперечные поляризации.
У длинноволновых акустических фононов $\omega=c_sk$, где $c_s$ --- скорость звука
соответствующей ветви; имеются две поперечные и одна продольная ветви.
В решётке число мод конечно, а их спектр ограничен атомным масштабом.

\subsection{Классический предел и закон Дюлонга--Пти}
При $T\gg\hbar\omega$ средняя энергия одной моды стремится к $T$:
по $T/2$ приходится на кинетическую и потенциальную части.
Если $T\gg\hbar\omega_{\max}$, все моды возбуждены классически и
\begin{equation*}
 E_{\mathrm{vib}}\simeq MT\simeq3NT,\qquad C_V\simeq M\simeq3N.
\end{equation*}
Это закон Дюлонга--Пти. В обычных единицах
$C_{V,\mathrm K}=k_BC_V\simeq3Nk_B$; для одного моля --- $3R$.

\clearpage
\section{Плотность состояний в модели Дебая}
\subsection{Подсчёт акустических мод}
В трёхмерном пространстве волновых векторов на одну ветвь приходится
\begin{equation*}
 d\mathcal N_s=\frac{V\,d^3k}{(2\pi)^3}
 =\frac{V}{2\pi^2}k^2\,dk
 =\frac{V\omega^2}{2\pi^2c_s^3}\,d\omega.
\end{equation*}
Суммируя две поперечные ветви со скоростью $c_t$ и одну продольную
со скоростью $c_l$, получаем
\begin{equation*}
 D(\omega)=\frac{V\omega^2}{2\pi^2}
 \left(\frac2{c_t^3}+\frac1{c_l^3}\right)
 =\frac{3V\omega^2}{2\pi^2c_D^3},\qquad
 \frac3{c_D^3}=\frac2{c_t^3}+\frac1{c_l^3}.
\end{equation*}
Здесь $D(\omega)d\omega$ --- число нормальных мод во всём образце
в интервале частот $d\omega$, а $c_D$ --- средняя дебаевская скорость звука.

\subsection{Предельная частота и температура Дебая}
Линейный акустический спектр продолжается до частоты $\omega_D$,
выбранной из условия сохранения полного числа мод:
\begin{align*}
 3N&=\int_0^{\omega_D}D(\omega)\,d\omega
 =\frac{V\omega_D^3}{2\pi^2c_D^3},\\
 \omega_D&=c_D\left(\frac{6\pi^2N}{V}\right)^{1/3},\qquad
 D_D(\omega)=
 \begin{cases}
  \dfrac{9N\omega^2}{\omega_D^3},&0<\omega<\omega_D,\\
  0,&\omega>\omega_D.
 \end{cases}
\end{align*}
Определим дебаевскую температуру в энергетических единицах
$\Theta_D=\hbar\omega_D$; её значение в кельвинах равно $\Theta_D/k_B$.

\begin{rbox}[unbreakable,left=14pt,right=14pt,top=10pt,bottom=8pt]{[]}
\centering
\captionsetup{font=footnotesize,hypcap=false,skip=5pt}
\begin{tikzpicture}[>=Stealth,font=\small,x=7cm,y=2.0cm]
 \draw[->] (0,0)--(1.55,0) node[right] {$\omega/\omega_D$};
 \draw[->] (0,0)--(0,1.7) node[above] {$D(\omega)$};
 \draw[p3,thick,domain=0:1,samples=60] plot (\x,\x^2);
 \draw[p3,thick] (1,1)--(1,0)--(1.45,0);
 \draw[dotted] (1,0)--(1,1.55);\node[below] at (1,0) {$1$};
 \draw[red3,dashed,thick] (0,0) .. controls (0.2,0) and (0.35,0.18) .. (0.5,0.47)
  .. controls (0.57,0.9) and (0.6,0.5) .. (0.66,0.62)
  .. controls (0.74,0.75) and (0.83,1.5) .. (0.88,1.4)
  .. controls (0.91,0.25) and (1,0.08) .. (1.25,0.02);
 \node[p3,right] at (0.18,1.25) {модель Дебая};
 \node[red3,right] at (1.04,0.85) {реальный спектр};
\end{tikzpicture}
\captionof{figure}{Дебаевская плотность состояний и схема более сложного реального спектра.
 Длинноволновая зависимость $D(\omega)\propto\omega^2$ сохраняется при малых частотах.}
\end{rbox}

\clearpage
\section{Энергия решётки и теплоёмкость}
\subsection{Нулевые и тепловые колебания}
Колебательная энергия образца равна сумме средних энергий всех мод:
\begin{align*}
 E_{\mathrm{vib}}&=\int_0^{\omega_D}D_D(\omega)
 \left[\frac{\hbar\omega}{2}
 +\frac{\hbar\omega}{e^{\hbar\omega/T}-1}\right]d\omega\\
 &=E_0+E_{\mathrm{th}}.
\end{align*}
Энергия нулевых колебаний в модели Дебая составляет
\begin{equation*}
 E_0=\frac{9N\hbar}{2\omega_D^3}\int_0^{\omega_D}\omega^3\,d\omega
 =\frac98N\hbar\omega_D=\frac98N\Theta_D.
\end{equation*}
Для тепловой части введём $z=\hbar\omega/T$:
\begin{align*}
 E_{\mathrm{th}}&=\frac{9N\hbar}{\omega_D^3}
 \int_0^{\omega_D}\frac{\omega^3\,d\omega}{e^{\hbar\omega/T}-1}\\
 &=\frac{9NT^4}{\Theta_D^3}
 \int_0^{\Theta_D/T}\frac{z^3\,dz}{e^z-1}.
\end{align*}
Полная энергия твёрдого тела дополнительно содержит статическую энергию
равновесной решётки: $E=U_{\mathrm{stat}}+E_0+E_{\mathrm{th}}$.
В гармоническом приближении при постоянном объёме частоты мод фиксированы;
$U_{\mathrm{stat}}$ и $E_0$ не зависят от температуры.

\subsection{Дебаевская формула для теплоёмкости}
Следовательно, $C_V=(\partial E_{\mathrm{th}}/\partial T)_V$.
Дифференцируя среднее заполнение каждой моды, получаем
\begin{align*}
 C_V&=\int_0^{\omega_D}D_D(\omega)
 \left(\frac{\hbar\omega}{T}\right)^2
 \frac{e^{\hbar\omega/T}}{(e^{\hbar\omega/T}-1)^2}\,d\omega\\
 &=9N\left(\frac{T}{\Theta_D}\right)^3
 \int_0^{\Theta_D/T}\frac{z^4e^z}{(e^z-1)^2}\,dz.
\end{align*}
Формула описывает колебательный вклад. В обычных единицах
$C_{V,\mathrm K}=k_BC_V$.

\clearpage
\section{Низкие и высокие температуры}
\subsection{Низкотемпературный закон Дебая}
При $T\ll\Theta_D$ верхний предел интеграла можно заменить бесконечностью:
\begin{align*}
 E_{\mathrm{th}}&\simeq\frac{9NT^4}{\Theta_D^3}
 \int_0^\infty\frac{z^3\,dz}{e^z-1}
 =\frac{9NT^4}{\Theta_D^3}\Gamma(4)\zeta(4)\\
 &=\frac{3\pi^4}{5}\frac{NT^4}{\Theta_D^3},\\
 C_V&\simeq\frac{12\pi^4}{5}N\left(\frac{T}{\Theta_D}\right)^3.
\end{align*}
Таким образом, колебательная теплоёмкость при низкой температуре
пропорциональна $T^3$. Возбуждаются главным образом длинноволновые
акустические моды; их число растёт как объём области $\hbar c_sk\lesssim T$
в пространстве волновых векторов.

\subsection{Высокотемпературное разложение}
При $T\gg\Theta_D$ во всём интеграле $z\ll1$. Используем
\begin{equation*}
 \frac1{e^z-1}=\frac1z-\frac12+\frac z{12}-\frac{z^3}{720}+O(z^5).
\end{equation*}
Положив $b=\Theta_D/T$, находим
\begin{align*}
 E_{\mathrm{th}}
 &=\frac{9NT^4}{\Theta_D^3}\int_0^b
 \left(z^2-\frac{z^3}{2}+\frac{z^4}{12}-\frac{z^6}{720}+\cdots\right)dz\\
 &=\frac{9NT^4}{\Theta_D^3}
 \left(\frac{b^3}{3}-\frac{b^4}{8}+\frac{b^5}{60}-\frac{b^7}{5040}+\cdots\right)\\
 &=3NT\left[1-\frac38\frac{\Theta_D}{T}
 +\frac1{20}\left(\frac{\Theta_D}{T}\right)^2
 -\frac1{1680}\left(\frac{\Theta_D}{T}\right)^4+\cdots\right].
\end{align*}
Постоянный член $-9N\Theta_D/8$ компенсирует $E_0$, если вернуть
нулевую энергию. Для полной колебательной энергии и теплоёмкости
\begin{align*}
 E_{\mathrm{vib}}&=3NT+\frac{3N\Theta_D^2}{20T}
 -\frac{N\Theta_D^4}{560T^3}+\cdots,\\
 C_V&=3N\left[1-\frac1{20}\left(\frac{\Theta_D}{T}\right)^2
 +\frac1{560}\left(\frac{\Theta_D}{T}\right)^4+\cdots\right]
 \longrightarrow3N.
\end{align*}

\clearpage
\section{Средний квадрат смещения атомов}
\subsection{Кинетическая и потенциальная энергии моды}
Для гармонического осциллятора средние кинетическая и потенциальная энергии
равны при любой температуре:
\begin{equation*}
 \left\langle\frac{P^2}{2m}\right\rangle
 =\left\langle\frac{m\omega^2Q^2}{2}\right\rangle
 =\frac12\overline E_\omega.
\end{equation*}
Отсюда
\begin{equation*}
 \langle Q^2\rangle_\omega
 =\frac{\hbar}{m\omega}\left(\frac12+\frac1{e^{\hbar\omega/T}-1}\right).
\end{equation*}
Обозначим через $\langle r^2\rangle$ средний квадрат смещения одного атома
от положения равновесия, усреднённый по всем атомам:
\begin{equation*}
 \langle r^2\rangle\equiv\frac1N\sum_{i=1}^N\langle|\mathbf u_i|^2\rangle
 =\frac{\hbar}{Nm}\int_0^{\omega_D}\frac{D_D(\omega)}{\omega}
 \left(\frac12+\frac1{e^{\hbar\omega/T}-1}\right)d\omega.
\end{equation*}
Ортогональность нормальных координат сохраняет сумму квадратов смещений.
Поэтому в модели Дебая
\begin{equation*}
 \langle r^2\rangle=\frac{9\hbar}{m\omega_D^3}
 \int_0^{\omega_D}\omega\left(\frac12+\frac1{e^{\hbar\omega/T}-1}\right)d\omega.
\end{equation*}

\subsection{Низкотемпературный предел}
После замены $z=\hbar\omega/T$ получаем
\begin{equation*}
 \langle r^2\rangle=\frac{9\hbar}{m\omega_D^3}
 \left[\frac{\omega_D^2}{4}
 +\left(\frac{T}{\hbar}\right)^2
 \int_0^{\Theta_D/T}\frac{z\,dz}{e^z-1}\right].
\end{equation*}
При $T\ll\Theta_D$ используем $\Gamma(2)\zeta(2)=\pi^2/6$:
\begin{equation*}
 \langle r^2\rangle\simeq\frac{9\hbar}{m\omega_D^3}
 \left[\frac{\omega_D^2}{4}+\frac{\pi^2}{6}\left(\frac{T}{\hbar}\right)^2\right].
\end{equation*}
При $T=0$ сохраняются нулевые смещения:
$\langle r^2\rangle_0=9\hbar/(4m\omega_D)$.

\subsection{Высокотемпературный предел}
При $T\gg\Theta_D$ слагаемые $1/2$ в нулевой и тепловой частях сокращаются:
\begin{equation*}
 \langle r^2\rangle\simeq\frac{9\hbar}{m\omega_D^3}
 \left(\frac{T\omega_D}{\hbar}+\frac{\hbar\omega_D^3}{36T}+\cdots\right)
 =\frac{9T}{m\omega_D^2}+\frac{\hbar^2}{4mT}+\cdots.
\end{equation*}

\clearpage
\section{Устойчивость решётки и соответственные состояния}
\subsection{Критерий Линдемана}
Эмпирический критерий связывает плавление с достижением определённой
относительной амплитуды колебаний:
\begin{equation*}
 \frac{\sqrt{\langle r^2\rangle_{T_m}}}{a}\simeq\xi,
 \qquad \xi\sim0.15\text{--}0.30,
 \qquad \langle r^2\rangle_{T_m}\simeq\xi^2a^2.
\end{equation*}
Здесь $a$ --- характерное межатомное расстояние, $T_m$ --- температура плавления.
Величина $\xi$ является ориентировочным параметром и зависит от вещества
и принятого определения амплитуды.
Если применим классический предел, то
\begin{equation*}
 \frac{9T_m}{m\omega_D^2}\simeq\xi^2a^2,
 \qquad T_m\simeq\frac{\xi^2m\omega_D^2a^2}{9}.
\end{equation*}
Для температуры плавления в кельвинах правую часть делят на $k_B$.

\subsection{Закон соответственных состояний в модели Дебая}
После введения отношения $y=T/\Theta_D$ температурные зависимости
колебательных величин принимают универсальный вид.
Например,
\begin{equation*}
 \frac{C_V}{N}=9y^3\int_0^{1/y}\frac{z^4e^z}{(e^z-1)^2}\,dz
 \equiv\mathcal C(y).
\end{equation*}
Для свободной энергии при фиксированном объёме
\begin{align*}
 F(V,T)&=F_0(V)+T\int_0^{\omega_D}D_D(\omega)
 \ln(1-e^{-\hbar\omega/T})\,d\omega\\
 &=F_0(V)+N\Theta_D(V)\,f\!\left(\frac{T}{\Theta_D(V)}\right),\\
 f(y)&=9y^4\int_0^{1/y}z^2\ln(1-e^{-z})\,dz.
\end{align*}
Величина $F_0(V)=U_{\mathrm{stat}}(V)+E_0(V)$ содержит статическую
и нулевую колебательную энергии.
Различия между веществами в этой модели входят через $N$,
фоновую энергию и масштаб $\Theta_D$.

\end{document}
